% Section commune : gaps FPGA (repris de gaps.md, 8e passage)
% Inclus par fpga.tex ET vision-embarquee.tex : pas de \ref nu vers sec:surete (absente de
% vision-embarquee.tex) ; le renvoi est protégé par \ifcsname r@sec:surete.
\section{Zones non couvertes et pistes de recherche}
\label{sec:gaps}

Les gaps ci-dessous viennent de \texttt{gaps.md} (8\textsuperscript{e} passage,
01/10/2026), établi sur les 101 fiches de la base et les 299 lignes de
\texttt{benchmarks.csv}, dont 216 sur FPGA. La carte compte 36 gaps, dont 29 étayés. Un
gap \emph{étayé} est constaté dans plusieurs fiches ; cela ne prouve pas qu'il est absent
de la littérature. Un gap à l'état d'\emph{hypothèse} ne repose que sur l'absence de
papiers dans la base. Le tableau~\ref{tab:gaps-fpga} reprend les gaps d'architecture, de
mesure et d'apprentissage ; le tableau~\ref{tab:gaps-surete} ceux de sûreté et de détection
de personnes, deux axes ajoutés à la base avec le thème de thèse%
\ifcsname r@sec:surete\endcsname{} (détail en section~\ref{sec:surete})\fi.

Trois gaps du passage précédent ont changé, parce qu'une fiche nouvelle les contredit :
\begin{itemize}
  \item \enquote{\glslink{ae}{autoencodeur} complet sur une petite FPGA} est supprimé. Un \glslink{cae}{CAE} complet,
    score par erreur de reconstruction, tourne en INT8 sur le \glslink{dpu}{DPU} d'une \glslink{zcu104}{ZCU104}
    \citep[§IV-B, Table~VI]{arxiv-2407.17647-energy-efficient-artefact-detection-accelerator},
    et une compression apprise place encodeur et décodeur sur \glslink{zcu102}{ZCU102}
    \citep[Table~X]{arxiv-2503.04832-lightweight-embedded-fpga-deployment-learned}\pp. Le
    coût des déconvolutions sur une puce plus petite passe dans le gap de l'anomalie
    visuelle ;
  \item \enquote{les apprentissages sur FPGA évitent la rétropropagation} devient
    \enquote{rétropropagation lente, en flottant, énergie non mesurée}. \glslink{ef-train}{EF-Train} fait
    FP\,/\,BP\,/\,WU en FP32 sur \glslink{zcu102}{ZCU102}, à 2,08\,\glslink{u-ms}{ms} par image pour un petit CNN
    CIFAR-10 à batch~128
    \citep[Table~7]{arxiv-2202.10935-ef-train-enable-efficient-device}\pp ; \glslink{chameleon}{Chameleon} le
    fait en FP16, à 413\,\glslink{u-ms}{ms} et 1,22\,\glslink{u-j}{J} par image, contre 33\,\glslink{u-ms}{ms} sur \glslink{jetson-nano}{Jetson Nano}
    \citep[Table~VI]{2024-aggarwal-chameleon-dual-memory-replay-online} ;
  \item \enquote{CL sur FPGA : précision faible} devient \enquote{énergie de portée non
    précisée}. \glslink{slda}{SLDA} sur \glslink{dpu}{DPU} atteint 78,79\,\% sur CORe50 (objet), à 9,00\,\glslink{u-w}{W} lus sur les
    rails sans portée déclarée
    \citep[Table~V]{2021-piyasena-accelerating-continual-learning-edge-fpga}.
\end{itemize}

\begin{table}[!htbp]
\centering\footnotesize
\caption{Gaps FPGA d'architecture, de mesure et d'apprentissage : statut dans le corpus et
piste principale.}
\label{tab:gaps-fpga}
\begin{tabularx}{\textwidth}{@{}P{4.6cm}P{2.5cm}L@{}}
\toprule
Gap & Statut & Piste \\
\midrule
\multicolumn{3}{@{}l}{\textit{Autoencodeurs}} \\
Anomalie visuelle (images RGB) par \glslink{ae}{AE} sur FPGA embarquée & étayé (15 fiches) ; un
  contre-exemple partiel (hyperspectral, seuil réglé sur le test) & chaîne \glslink{cae}{CAE} INT8 de
  Vitis~AI sur MVTec AD et \glslink{kv260}{KV260}, seuil fixé hors du test, \glslink{auroc}{AUROC} et puissance de carte ;
  coût des déconvolutions en dataflow \\
Compression apprise embarquée avec codage entropique & hypothèse (4 fiches) & chaîne
  complète sur plateforme à basse consommation, énergie par image mesurée \\
\glslink{vae}{VAE} en matériel pour un latent de grande dimension & hypothèse (2 fiches) & coût d'un
  générateur gaussien parallèle ; score $D_{KL}$ déterministe en alternative \\
\midrule
\multicolumn{3}{@{}l}{\textit{Détection}} \\
Protocole de mesure commun pour \glslink{yolo}{YOLO} sur FPGA & étayé (24 fiches) & \glslink{map}{mAP} sur la cible,
  batch~1 de bout en bout avec pire cas et gigue, puissance de puce et de carte \\
Détection en dataflow quantifié sur SoC à mémoire limitée & étayé (13 fiches), nuancé &
  \glslink{yolo}{YOLO} récent en W4--W8 par FINN ou SATAY sur \glslink{kv260}{KV260}, face à la \glslink{dpu}{DPU} sur la même carte \\
Petits objets : effet de la quantification par taille d'objet & étayé (2 fiches) pour
  l'absence ; hypothèse sinon & \glslink{ap}{AP}\textsubscript{small} COCO mesurée sur la cible \\
Post-traitement et interaction PS--PL & étayé (14 fiches) & \glslink{nms}{NMS} matériel à côté de la \glslink{dpu}{DPU},
  latence de bout en bout et borne \\
\midrule
\multicolumn{3}{@{}l}{\textit{Mesure et comparaison}} \\
Précision non vérifiée sur la cible & étayé (20 fiches), nuancé & \glslink{map}{mAP} COCO val sur la
  carte, étape par étape (quantification, compilation, opérateurs, \glslink{nms}{NMS}) \\
Baseline \glslink{jetson}{Jetson} récente mesurée dans les mêmes conditions & étayé (27 fiches) & même modèle,
  même précision numérique, même batch, même méthode de puissance sur \glslink{kv260}{KV260} et \glslink{orin-nano}{Orin Nano} \\
Puissance puce contre carte & étayé (40 fiches + \texttt{benchmarks.csv}) & télémétrie de
  la puce + wattmètre de carte + puissance au repos sur la même carte \\
Métriques flatteuses et comparaisons biaisées & étayé (20 fiches) & recalculer les rapports
  annoncés à métriques imposées \\
Code et bitstreams rarement publiés & étayé ($\approx$\,55 fiches expérimentales) & publier
  projet, bitstream et scripts de mesure \\
\midrule
\multicolumn{3}{@{}l}{\textit{Toolflows}} \\
Comparaison croisée hls4ml / FINN / Vitis AI & étayé (12 fiches), nuancé & un même \glslink{yolo}{YOLO} à
  opérateurs \glslink{relu}{ReLU} exporté en QONNX, trois flots, une carte \\
Réseaux qui ne tiennent pas on-chip & étayé (13 fiches) & quantification agressive + FCMP
  ou éviction SMOF, coût à batch~1 et variabilité temporelle \\
Modèle de puissance validé dans les toolflows & hypothèse (outils) ; étayé (7 fiches) pour
  l'estimation non validée & vérifier la documentation actuelle de FINN et hls4ml \\
NAS et quantification mixte guidées par des mesures FPGA & hypothèse (nuancée) & HAWQ + ILP
  appliqués à la tête d'un \glslink{yolo}{YOLO} \\
\midrule
\multicolumn{3}{@{}l}{\textit{Apprentissage continu}} \\
Rétropropagation sur SoC FPGA : lente, en flottant, énergie non mesurée & étayé (7 fiches)
  & queue de réseau en précision réduite en PL, étage figé sur la \glslink{dpu}{DPU} \\
Adaptation à la dérive en détection d'anomalies & hypothèse (5 fiches) & \glslink{onlad}{ONLAD} sur les
  caractéristiques d'un encodeur CNN figé \\
CL sur FPGA : énergie de portée non précisée ou estimée & étayé (6 fiches) & remesurer \glslink{slda}{SLDA}
  avec wattmètre et télémétrie \\
Apprentissage continu limité à la classification & étayé (13 fiches), nuancé (détection
  incrémentale sans FPGA) & Latent Distillation sur une tête \glslink{yolo}{YOLO}, backbone figé sur la \glslink{dpu}{DPU}
  \\
\bottomrule
\end{tabularx}
\end{table}

\begin{table}[!htbp]
\centering\footnotesize
\caption{Gaps de sûreté, de certification et de détection de personnes : statut dans le
corpus et piste principale. Aucun texte normatif n'est dans la base.}
\label{tab:gaps-surete}
\begin{tabularx}{\textwidth}{@{}P{4.6cm}P{2.5cm}L@{}}
\toprule
Gap & Statut & Piste \\
\midrule
\multicolumn{3}{@{}l}{\textit{Sûreté, fiabilité et certification}} \\
Aucun détecteur quantifié sur FPGA dans les travaux de fiabilité, de vérification ou de
  surveillance & étayé (25 fiches) & un \glslink{yolov3-tiny}{YOLOv3-tiny} ou v8n sur \glslink{kv260}{KV260}, objet commun aux
  expériences de fiabilité, de \glslink{map}{mAP} sur cible et de moniteur \\
Fiabilité SEU d'un détecteur : criticité par couche, TMR sélectif, \glslink{dpu}{DPU} contre dataflow &
  étayé (9 fiches) & injection (configuration, \glslink{bram}{BRAM}, multi-fautes) sur un \glslink{yolo}{YOLO} INT8 FINN et
  \glslink{dpu}{DPU} ; durcissement sélectif chiffré en \glslink{lut}{LUT} et \glslink{bram}{BRAM} \\
Critère d'erreur dangereuse pour la détection de personnes ; \glslink{fit}{FIT} et précision vers \glslink{sil}{SIL} / PL
  & étayé (11 fiches) & non-détection définie sur une séquence ; taux de défaillance par
  heure combinant inexactitude et fautes \\
Preuve d'équivalence modèle $\rightarrow$ bitstream ; qualification des toolflows & étayé
  (14 fiches) & sorties bit à bit du modèle QONNX et du bitstream ; analyse TQL de FINN ou
  Vitis~AI \\
Vérification formelle non scalable aux CNN de détection & étayé (4 fiches) ; hypothèse pour
  les outils récents & vérifier un moniteur ou un post-traitement RTL, plus petit que le
  réseau \\
Moniteurs d'exécution (OOD, incertitude, conforme) non portés sur FPGA & étayé (20 fiches)
  & prédiction conforme, puis moniteur à boîtes en PL ; surcoût mesuré sur \glslink{kv260}{KV260} \\
Correspondance des cadres aéronautique et automobile avec les normes industrielles & étayé
  (11 fiches) ; hypothèse pour la littérature & tableau objectif par objectif vers IEC 61508,
  ISO 13849 et ISO/IEC TR 5469 \\
Déterminisme du FPGA : \glslink{wcet}{WCET} et gigue non mesurés & étayé (12 fiches) & traces de latence
  sur $10^5$ images, \glslink{dpu}{DPU}, FINN et TensorRT \\
Apprentissage en ligne exclu des cadres de certification & étayé (13 fiches) & \glslink{slda}{SLDA}
  StaticCov avec test de non-régression avant activation \\
\midrule
\multicolumn{3}{@{}l}{\textit{Détection de personnes en milieu industriel}} \\
Jeu de données et cas d'assurance en zone industrielle & étayé (8 fiches) ; hypothèse pour
  la littérature & recherche de jeux EPI ; ODD d'exemple selon AMLAS \\
Détecteurs de personnes quantifiés sur FPGA & étayé (9 fiches) ; un seul travail mesuré &
  \glslink{yolov8n}{YOLOv8n} sur CrowdHuman, rappel à \glslink{fppi}{FPPI} fixé sur \glslink{kv260}{KV260}, flottant et INT8 \\
Chaîne détection $\rightarrow$ suivi $\rightarrow$ ré-identification sur FPGA & étayé (4
  fiches) ; hypothèse pour la Re-ID & suivi simple (SORT) sur le PS de la \glslink{kv260}{KV260} \\
Robustesse adverse des détecteurs quantifiés & hypothèse (3 fiches) & rejouer un patch
  adverse sur le \glslink{yolo}{YOLO} quantifié, avec et sans moniteur \\
\bottomrule
\end{tabularx}
\end{table}

\gap{la sûreté du détecteur embarqué} Les deux moitiés du sujet de thèse existent, mais
séparément. Les études de fiabilité sur FPGA portent sur de petits classifieurs : 19 et
1\,848 poids
\citep[§I, §III]{2019-libano-selective-hardening-neural-networks-fpgas}, un \glslink{qnn}{QNN} FINN sur
MNIST et CIFAR-10
\citep[§V]{arxiv-1912.07394-efficient-error-tolerant-quantized-neural}. Inversement,
aucun des 18 papiers de l'axe \texttt{yolo-fpga} ne traite la sûreté, et le seul \glslink{yolo}{YOLO}
\enquote{certifié} du corpus tourne sur \glslink{jetson}{Jetson}
\citep[§V]{arxiv-2406.12362-certified-ml-object-detection-surveillance}. Les guides
renvoient l'implémentation à plus tard : l'EASA écarte la conformité du matériel
spécialisé \citep[§4.2.3]{2020-easa-concepts-design-assurance-neural-networks} et
n'accepte pas l'apprentissage en ligne
\citep[MOC LM-02]{2024-easa-concept-paper-guidance-level-1}.

\gap{la détection de personnes sur FPGA} Un seul papier mesure un détecteur de personnes
sur FPGA : \glslink{yolov8}{YOLOv8}, v11 et v26 nano en INT8 sur \glslink{zcu104}{ZCU104}, \glslink{map5095}{mAP@0.5:0.95} de 0,66 à 0,79 sur un
jeu maison non publié, 34,05\,\glslink{fps}{FPS} de bout en bout et une puissance estimée par Vivado
\citep[Table~II, §IV-C]{arxiv-2607.13106-no-attention-no-problem-dpu}\pp. Aucun travail
ne donne de taux de personnes manquées à taux de fausse alarme fixé.

\subsection{Quatre questions de thèse candidates}

Les pistes suivantes sont des propositions de ce document \deduit. Elles combinent des gaps
étayés et s'appuient sur des briques déjà présentes dans le corpus. Elles partagent un
même objet expérimental : un \glslink{yolo}{YOLO} quantifié de détection de personnes sur \glslink{kv260}{KV260}.

\begin{enumerate}
  \item \textbf{Un détecteur de personnes certifiable sur SoC FPGA.} Établir, pour un même
    \glslink{yolo}{YOLO}, une chaîne de preuve d'équivalence entre le modèle QONNX et le bitstream, une borne
    de latence mesurée (pire cas, gigue) et un moniteur d'exécution sur la même puce.
    Briques existantes :
    \begin{itemize}
      \item le \glslink{yolov3-tiny}{YOLOv3-tiny} allégé en w2a4 par FINN sur \glslink{xc7z020}{Zynq-7020}
        \citep{arxiv-2609.24757-npu-accelerator-quantized-real-time}\pp, flot dataflow plus
        inspectable que la \glslink{dpu}{DPU} \deduit ;
      \item la criticité par couche et le TMR sélectif de
        \citet{2019-libano-selective-hardening-neural-networks-fpgas} ;
      \item la prédiction conforme sur \glslink{yolo}{YOLO}
        \citep{arxiv-2304.06052-confident-object-detection-conformal-prediction}\pp{} et le
        moniteur à boîtes d'abstraction
        \citep{arxiv-1911.09032-outside-box-abstraction-monitoring-neural}.
    \end{itemize}
    Verrou : relier ces mesures à un taux de défaillance dangereuse par heure, dans un cadre
    de sécurité machine que le corpus ne contient pas encore.
  \item \textbf{Un protocole de mesure et une comparaison croisée des flots.} Un même \glslink{yolo}{YOLO}
    compilé par Vitis~AI, et par FINN ou hls4ml via QONNX, sur \glslink{kv260}{KV260}, face à une \glslink{orin}{Jetson Orin}
    Nano. Mesures : \glslink{map}{mAP} COCO et rappel de personnes sur la cible, latence de bout en bout et
    sa distribution, puissance de puce et de carte. Le protocole partirait du cadre
    d'\citet{arxiv-2304.11580-framework-benchmarking-real-time-embedded} (\glslink{map}{mAP} sur la carte,
    puissance de carte, repos déduit) et de ceux de \glslink{dac-sdc}{DAC-SDC}
    \citep{2021-xu-dac-sdc-low-power-object}, de FINN
    \citep{arxiv-1612.07119-finn-framework-fast-scalable-binarized} et d'Alqahtani et al.
    \citep{arxiv-2409.16808-comprehensive-evaluation-deep-learning-object}\pp.
  \item \textbf{Détection d'anomalies visuelles par \glslink{ae}{autoencodeur} convolutif sur SoC FPGA.}
    Briques existantes :
    \begin{itemize}
      \item le \glslink{cae}{CAE} INT8 sur \glslink{dpu}{DPU} et son score par reconstruction
        \citep{arxiv-2407.17647-energy-efficient-artefact-detection-accelerator}, dont le
        seuil est réglé sur le test ;
      \item l'encodeur ternaire paramétrable de
        \citet{2021-manjunath-energy-efficient-edgeai-autoencoder-accelerator} ;
      \item le score calculé sur la puce et le \glslink{vae}{VAE} réduit à l'encodeur de
        \citet{arxiv-2108.03986-autoencoders-fpgas-real-time-unsupervised}.
    \end{itemize}
    Verrou : une \glslink{auroc}{AUROC} mesurée sur la carte, sur MVTec AD, avec un seuil fixé hors du test,
    sur \glslink{kv260}{KV260} ou en dataflow. Le même score pourrait servir de moniteur hors distribution
    pour le détecteur \deduit.
  \item \textbf{Adaptation en ligne certifiable par construction.} Un étage figé INT8 sur
    la \glslink{dpu}{DPU} et une petite tête apprenante : \glslink{slda}{SLDA}
    \citep{2021-piyasena-accelerating-continual-learning-edge-fpga}, recette \glslink{qlr-cl}{QLR-CL}
    \citep{arxiv-2110.10486-tinyml-platform-device-continual-learning} ou \glslink{os-elm}{OS-ELM} pour la
    dérive \citep{2019-tsukada-neural-network-device-learning-anomaly}. Un test de
    non-régression sur un jeu de référence embarqué et un retour arrière précéderaient
    toute activation des nouveaux paramètres. Mesures : latence et énergie par mise à jour,
    surcoût de la validation. L'extension à la détection reste ouverte : la détection
    incrémentale n'existe que sans FPGA
    \citep{arxiv-2409.01872-latent-distillation-continual-object-detection}\pp.
\end{enumerate}
